% Part: first-order-logic
% Chapter: beyond
% Section: introduction

\documentclass[../../../include/open-logic-section]{subfiles}

\begin{document}

\olfileid{fol}{byd}{int}

\olsection{In grote lijnen}

De eerste-ordelogica is niet de enige logica die de moeite waard is.
Er bestaan veel uitbreidingen en varianten. Een logica geeft doorgaans
precies aan welke formele taal zij gebruikt. Daar horen meestal, maar
niet altijd, regels bij waarmee je conclusies mag afleiden (een
deductiesysteem). Meestal wordt ook uitgelegd wat de tekens bedoeld zijn
te betekenen (de bedoelde semantiek), maar ook dat hoeft niet. Het
technische woord ``logica'' roept dan een vraag op. Wat hebben systemen
die geen eerste-ordelogica zijn te maken met ``logica'' in de gewone,
intuïtieve of filosofische betekenis? Elk systeem hieronder is gemaakt
om een bepaalde manier van redeneren weer te geven. Maar kunnen we ook
zeggen wat zo'n systeem logisch maakt?

Daar is geen eenvoudig antwoord op. Mensen gebruiken het woord
``logica'' op verschillende manieren en in verschillende situaties.
Filosofen hebben ook heel verschillend uitgelegd wat ``logica'' en
``waarheid'' zijn. Volgens één opvatting moet logisch redeneren uitmaken
welke uitspraken noodzakelijk waar zijn. Volgens een andere opvatting
gaat het om uitspraken waarvan we de waarheid zonder zintuiglijke
waarneming kunnen vaststellen: uitspraken die a priori waar zijn. Weer
andere opvattingen kijken naar uitspraken die waar blijven hoe je de
niet-logische termen ook uitlegt, die waar zijn door hun logische vorm,
of die waar zijn doordat we taalafspraken hebben gemaakt. Elk van deze
vijf ideeën moet nog veel preciezer worden uitgelegd. Zelfs onder
mensen die alleen kijken naar de logica die wiskundigen gebruiken, is
er weinig overeenstemming over wat daarbij hoort. Russell en Whitehead
probeerden bijvoorbeeld in de \textit{Principia Mathematica} de
wiskunde helemaal op logica te bouwen. Daarmee volgden ze de door Frege
begonnen stroming die {\em logicisme} heet. Ze gebruikten een logica met
hogere typen, die lijkt op een systeem dat hieronder wordt beschreven.
Uiteindelijk moesten ze wel axioma's toevoegen die volgens de meeste
maatstaven niet zuiver logisch zijn. Het ging vooral om het
oneindigheidsaxioma en het axioma van reducibiliteit. Dat laatste zegt,
grofweg, dat een eigenschap van hogere orde door een eigenschap van de
eerste orde kan worden vervangen. Toch kun je vinden dat sommige
manieren van redeneren over eigenschappen of andere hogere-ordeobjecten
wel degelijk logisch zijn. Quine denkt daar anders over. In zijn
ontologie---zijn opvatting over welke soorten dingen als object mogen
meetellen---zijn ``proposities'' geen geldige objecten waarover je kunt
spreken. Volgens hem zijn tweede-ordelogica en nog hogere-ordelogica
eigenlijk verzamelingenleer in vermomming. Zijn punt is dus dat een
systeem niet zuiver logisch is zodra daarin bijvoorbeeld ``voor elke
eigenschap'' staat en het systeem daarmee over predicaten
kwantificeert.

Voor nu laten we deze filosofische vragen maar liggen. We behandelen de
systemen hieronder gewoon als formele idealiseringen: nauwkeurig
gemaakte versies van verschillende manieren van redeneren, of die nu
logisch heten of niet.

\end{document}
